\section{Results}
\label{S:results}

The principal SAL tables are recomputed from saved row-level artifacts. The other tables distinguish
aggregate result files from development notes. All values describe the recorded configurations;
none is treated as an untouched-test estimate of a universal final pipeline.

\subsection{RQ1: localization, detection, and verification}

\paragraph{Localizer development.}
The internal ten-paper notebook retains executed outputs for the terminal-window and
mid-paper-anchor localizers. The former keeps 433 sentences and the latter 559, while reported
localization at 0.70 token containment rises from 0.633 to 0.857. Drop rates are 87.3\% and
83.6\%. This reference includes scope and assumption statements as well as limitations and
future work; these are region-coverage observations, not strict two-class extraction scores.
At 0.90 containment, localization is lower (0.429 and 0.633), showing dependence on the overlap
criterion. Full values and the logged section-aware comparison are in \Cref{A:development}.

The separate SAL-training development notebook reports five grouped-fold summaries of 0.861
mean reachability for anchor localization and 0.988 for whole target sections. The latter uses
supplied section names; it does not measure PDF heading recovery. Its cue-augmented count adds
potentially overlapping windows without full deduplication, so the archived augmented drop rates
are approximate. These summaries describe a development protocol with token-containment scoring;
they are not alternative estimates of the strict 40-paper evaluation below.

\paragraph{Fixed SAL test stream.}
Of 7,968 SAL sentences, 1,885 belong to the recovered Discussion or Conclusion sections.
This gate discards 76.3\% of the input while retaining 203 of 205 labelled limitations
(reachability 0.990). Prevalence rises from 2.57\% in the full stream to 10.77\% in the retained
pool. This is a conditional localization result using recovered tags rather than automatically
parsed PDF headings.

\begin{table}[t]
\centering
\small
\caption{SAL stage contributions on 40 articles, with all 205 labelled positives in the recall
denominator. TP denotes retained labelled positives. Both Stage-C rows use complete cached
judgments on the same 391 PubMedBERT candidates.}
\label{T:sal-stages}
% Auto-generated from exact saved prediction/verdict IDs
\begin{tabular}{@{}lrrrrr@{}}
\toprule
Configuration & Kept & TP & Precision & Recall & $F_1$ \\
\midrule
Stage A & 1885 & 203 & 0.108 & 0.990 & 0.194 \\
A+B: PubMedBERT & 391 & 180 & 0.460 & 0.878 & 0.604 \\
A+B+C: RCT prompt & 269 & 156 & 0.580 & 0.761 & 0.658 \\
A+B+C: SAL keep-biased & 249 & 148 & 0.594 & 0.722 & 0.652 \\
\bottomrule
\end{tabular}

\end{table}

\Cref{T:sal-stages} shows the resulting trade-off. PubMedBERT after localization retains
391 sentences, including 180 labelled positives. The RCT verifier removes 98 of the 211 detector
false positives (46.4\%), but also removes 24 of its 180 true positives. True-positive retention
is therefore 86.7\%. Precision increases from 0.460 to 0.580; recall decreases from 0.878 to 0.761;
$F_1$ increases from 0.604 to 0.658. The SAL keep-biased alternative retains fewer sentences
and achieves slightly higher precision, 0.594, but lower recall, 0.722, and $F_1$, 0.652.

The paired descriptive bootstrap interval for the RCT-prompt $F_1$ change is
$[0.026,0.083]$, around an observed increase of 0.054. The corresponding intervals for A+B and
A+B+C are $[0.554,0.654]$ and $[0.604,0.710]$. These intervals concern the saved-run comparison
only. Prompt refinement used benchmark feedback, so the positive interval does not establish a
confirmatory generalization claim.

An earlier notebook reports $F_1=0.701$ after judging a union of several detectors' candidates.
That run and the candidate-specific cached run are not interchangeable, because batch composition
and prompt configuration differ. We use the latter for the principal result. Nor is our full-stream
$F_1$ directly comparable to the published 0.821 recognition result of
\citet{lan2024automatic}, which used a different candidate-selection and model-development protocol.

\begin{table}[t]
\centering\small
\caption{SAL prevalence sensitivity with fixed predictions/verdicts. The section pool has
203 reachable positives; matching retains them and samples 778 negatives, averaging over eight
seeds. Full-stream scores retain all 205 positives in the recall denominator.}
\label{T:prevalence}
% Generated offline by revision_audit/verify_minimal.py
\begin{tabular}{@{}lrrrr@{}}
\toprule
Configuration & Full-stream $F_1$ & Pool recall & Pool $F_1$ & Matched $F_1$ \\
\midrule
A+B & 0.604 & 0.887 & 0.606 & 0.749 \\
A+B+C: RCT & 0.658 & 0.768 & 0.661 & 0.757 \\
A+B+C: SAL & 0.652 & 0.729 & 0.655 & 0.741 \\
\bottomrule
\end{tabular}

\end{table}

The recomputed prevalence analysis makes this distinction explicit (\Cref{T:prevalence}).
With negative subsampling to approximately 20.7\%, detector $F_1$ is 0.749, RCT-verified $F_1$
is 0.757, and SAL-verified $F_1$ is 0.741. Thus the RCT verifier's advantage narrows, and the SAL
variant no longer improves this measure. These are rescorings of fixed decisions, not new LLM
runs. They demonstrate dependence on the evaluated negative pool, without establishing a
universal prevalence threshold at which verification becomes useful.

Other prompt and model experiments have different judging populations. The stored
120-candidate balanced hard-sample sweep reports $F_1$ of 0.789, 0.779, and 0.784 for the three
Yandex variants. Aggregate union-candidate runs include the strict SAL result of 0.568 and
keep-biased SAL result of 0.554 for PubMedBERT; neither replaces its candidate-specific result
of 0.652. Model, prompt and batch membership differ across some runs, preventing attribution
of those differences solely to prompt wording. \Cref{A:verifiers} reports these development
results separately.

\subsection{RQ2: detector and reference dependence}

\begin{table}[t]
\centering
\small
\caption{Eight saved SAL detector vectors scored on the full stream and after the same
Discussion/Conclusion gate. These are positive-class metrics; the artifact permits metric
recomputation but does not fully reconstruct each training job.}
\label{T:sal-detectors}
% Auto-generated by revision_audit/verify_results.py
\begin{tabular}{@{}lrrrr@{}}
\toprule
Stage-B detector & Full-stream $F_1$ & A+B precision & A+B recall & A+B $F_1$ \\
\midrule
MiniLM-L6 (frozen) & 0.216 & 0.249 & 0.727 & 0.371 \\
mpnet (frozen) & 0.280 & 0.304 & 0.776 & 0.437 \\
bge-small (frozen) & 0.282 & 0.298 & 0.810 & 0.436 \\
bge-base (frozen) & 0.294 & 0.299 & 0.815 & 0.437 \\
e5-base (frozen) & 0.314 & 0.312 & 0.849 & 0.457 \\
DistilBERT-FT & 0.413 & 0.377 & 0.849 & 0.523 \\
SciBERT-FT & 0.458 & 0.405 & 0.873 & 0.553 \\
PubMedBERT-FT & 0.495 & 0.460 & 0.878 & 0.604 \\
\bottomrule
\end{tabular}

\end{table}

Localization raises $F_1$ for every saved SAL detector vector (\Cref{T:sal-detectors}).
PubMedBERT has the highest A+B $F_1$ in this comparison, followed by SciBERT and DistilBERT.
The frozen heads retain useful recall but produce more false positives. This supports the
localization mechanism for these fixed predictions; it does not establish that parameter count,
pretraining domain, or fine-tuning alone caused the observed ranking.

\paragraph{Classical and frozen baselines.}
The saved native-prevalence baseline grid is reported in \Cref{T:baseline-grid}. On the full
SAL stream, TF-IDF/logistic regression has $F_1=0.443$, compared with 0.273 for its separately
fitted bge-small head; the cached bge-small column in \Cref{T:sal-detectors} has 0.282. These are
different saved fits, not rounding variants of one run. On the grid's FWS split, TF-IDF/logistic
regression reaches 0.714, and on binary Sci-Sentence the in-domain bge-small head reaches 0.814.
All-negative predictions have zero positive-class $F_1$. These results establish useful
classical baselines and show why encoder size alone is an inadequate explanation of performance.
The FWS grid score must not be treated as a paired comparison with the 0.707 breadth result,
because their test pools differ.

\begin{table}[t]
\centering\small
\caption{Saved baseline-grid aggregates: positive-class metrics. SAL uses the 7,968-row full
stream; FWS uses a separate paper-level 20\% holdout, whose exact test count is not stored in
the aggregate; SciSentence uses the supplied 140-row binary test.}
\label{T:baseline-grid}
% Generated offline by revision_audit/verify_minimal.py
\begin{tabular}{@{}llrrr@{}}
\toprule
Subset & Method & Precision & Recall & $F_1$ \\
\midrule
SAL & majority (all-neg) & 0.000 & 0.000 & 0.000 \\
SAL & TF-IDF + LogReg & 0.362 & 0.571 & 0.443 \\
SAL & CountBin(1-4) + BernoulliNB & 0.000 & 0.000 & 0.000 \\
SAL & cue-rules only & 0.355 & 0.215 & 0.267 \\
SAL & bge-small frozen + logreg & 0.164 & 0.820 & 0.273 \\
FWS & majority (all-neg) & 0.000 & 0.000 & 0.000 \\
FWS & TF-IDF + LogReg & 0.667 & 0.767 & 0.714 \\
FWS & CountBin(1-4) + BernoulliNB & 0.880 & 0.510 & 0.646 \\
FWS & cue-rules only & 0.855 & 0.559 & 0.676 \\
FWS & bge-small frozen + logreg & 0.514 & 0.820 & 0.632 \\
SciSentence & majority (all-neg) & 0.000 & 0.000 & 0.000 \\
SciSentence & TF-IDF + LogReg & 0.784 & 0.725 & 0.753 \\
SciSentence & CountBin(1-4) + BernoulliNB & 0.895 & 0.425 & 0.576 \\
SciSentence & cue-rules only & 1.000 & 0.150 & 0.261 \\
SciSentence & bge-small frozen + logreg & 0.761 & 0.875 & 0.814 \\
\bottomrule
\end{tabular}

\end{table}

\begin{table}[t]
\centering
\small
\caption{Aggregate breadth results. FWS targets author-stated future work; Sci-Sentence merges
Limitation and Research Gap and includes prior-work criticism. Neither row evaluates raw-PDF
localization. The Sci-Sentence filter here uses the author-gap prompt and is deliberately
reported as a target mismatch.}
\label{T:breadth}
\begin{tabular}{@{}llrrr@{}}
\toprule
Subset & Configuration & Precision & Recall & $F_1$ \\
\midrule
FWS ($n=1{,}506$) & A+B & 0.723 & 0.618 & 0.667 \\
 & A+B+C, future-work prompt & 0.843 & 0.609 & 0.707 \\
Sci-Sentence ($n=140$) & Detector & 0.578 & 0.650 & 0.612 \\
 & Detector+C, author-gap prompt & 0.900 & 0.450 & 0.600 \\
\bottomrule
\end{tabular}
\end{table}

The aggregate FWS run raises precision with little recorded recall loss, whereas the
Sci-Sentence run raises precision at a larger recall cost and slightly reduces $F_1$
(\Cref{T:breadth}). The latter's label definition permits prior-work gaps that an author-gap
prompt rejects. Subsequent notes record target-matched prompt and in-domain threshold adaptation,
but complete decisions for those additional runs are not available in the supplied snapshot.
The development log reports $F_1=0.730$ for the intended gap-matched prompt with the existing
cross-domain detector, and 0.880 for an in-domain bge-small head at threshold 0.3 followed by
that prompt. Before verification, the latter retains 39 of 40 positives among 61 candidates;
after verification it retains 33 positives among 35 candidates. These logged counts are
arithmetically consistent with precision 0.943 and recall 0.825, but their full decisions are
not supplied. The contradictory shared prompt suffix described in \Cref{S:stage-c}, together
with test-informed development, limits interpretation as a clean policy-alignment experiment.
\Cref{A:breadth-adaptation} records the configurations; none is an independent confirmation
that the verifier is nearly error-free.

\begin{table}[t]
\centering
\small
\caption{Internal arXiv development notes, restricted to the 11 limitation and 14 future-work
reference statements. The 49-statement mixed-type denominator is not used. Retained-sentence
counts indicate candidate load, not validated precision. No new training replication was performed.}
\label{T:transfer}
\begin{tabular}{@{}lrrr@{}}
\toprule
Rules + detector & \shortstack{Limitation\\recall} & \shortstack{Future-work\\recall} & \shortstack{Retained\\sentences} \\
\midrule
bge-small (frozen; mixed-source training) & 0.545 & 0.786 & 97 \\
MiniLM-L6 (frozen; mixed-source training) & 0.636 & 0.857 & 170 \\
mpnet (frozen; mixed-source training) & 0.455 & 0.786 & 186 \\
PubMedBERT (fine-tuned; mixed-source training) & 0.364 & 0.714 & 65 \\
\bottomrule
\end{tabular}
\end{table}

The internal reference illustrates a different operating-point trade-off
(\Cref{T:transfer}). MiniLM retains more target statements and more candidates than bge-small.
The full fine-tuned model has lower recorded target recall on this small silver set. The result
cannot establish that frozen encoders generally transfer better: the set is small, its reference
coverage is incomplete, training-label mappings differ from the principal task, and precision
was not independently measured. It also explains why the selected development notebook and the
portable runner currently name different default heads. The complete logged head-only and
hybrid comparisons, class-balancing checks, and within-corpus benchmark tables are retained in
\Cref{A:development}. Validation macro-$F_1$ and positive-reference recovery answer different
questions. Lower recovery after transfer is compatible with distribution shift, label-policy
mismatch, incomplete silver references, or overfitting; these observations do not isolate one
cause or establish a universal frozen-versus-fine-tuned ranking.

Alternative references further change the interpretation. The development notes describe a
33-statement two-target model reference and a stricter 19-statement reference. MiniLM's reported
recovery changes from 0.455 to 0.632, while Stage-A reachability changes from 23/33 to 14/19.
The latter reference was collected with incomplete reading of three long papers. This is
sensitivity to annotation coverage and policy on the same articles, not a measured system
improvement. The larger documented 230-paper reference and its pending human protocol are
summarized in \Cref{A:reference}. They supply no completed human accuracy result.

Neighbor-position scoring on the saved SAL stream also changes the numbers. For A+B, $F_1$
rises from 0.604 under strict scoring to 0.743 at $W=1$; RCT-verified output rises from 0.658 to
0.821. The latter number is coincidentally equal, after rounding, to the published SAL detector
score under a different protocol. It is not evidence of matching that system. Neighbor scoring
can credit multiple predictions around one gold sentence and does not establish semantic
correctness or repaired annotations (\Cref{T:neighbor}).

\subsection{RQ3: input volume and full-article prompting}

\begin{table}[t]
\centering
\small
\caption{Estimated input volume over the 40 SAL articles. The accounting includes repeated
verification prompts and local context, using \texttt{cl100k\_base} and approximate wrapper
constants. It excludes output tokens, inference latency, local detector work, and retries.}
\label{T:tokens}
\begin{tabular}{@{}lrr@{}}
\toprule
Input accounting & Total input tokens & Mean per article \\
\midrule
Full-article prompting & 282,846 & 7,071.2 \\
Selective verification & 75,622 & 1,890.6 \\
\bottomrule
\end{tabular}
\end{table}

The accounting file gives a ratio of $282{,}846/75{,}622=3.740$,
a 73.3\% reduction in estimated LLM input (\Cref{T:tokens}). Only 391 of 7,968 source sentences
reach verification, about 9.8 candidates per article. The token ratio is smaller than the
sentence-count ratio because repeated prompts and neighboring context contribute substantial
overhead. It is not a measured monetary saving, throughput increase, or constant per-article cost.

\begin{table}[t]
\centering
\small
\caption{Recorded ten-article SAL comparison: 2,053 sentences and 55 positives, sampled from
articles having at least two labelled limitations. The complete stage-count file survives;
full-article extraction survives as aggregate precision/recall and mean API input only.
The baseline $F_1$ is approximately 0.427, computed from its rounded precision and recall.}
\label{T:full-paper}
\begin{tabular}{@{}lrrrr@{}}
\toprule
Configuration & Kept & Precision & Recall & $F_1$ \\
\midrule
Stage A & 498 & 0.110 & 1.000 & 0.199 \\
A+B: PubMedBERT & 109 & 0.468 & 0.927 & 0.622 \\
A+B+C: RCT verifier & 77 & 0.584 & 0.818 & 0.682 \\
Full-article LLM & --- & 0.559 & 0.345 & $\approx$0.427 \\
\bottomrule
\end{tabular}
\end{table}

In this selected sample, A+B+C has higher recorded recall and $F_1$ than the full-article
baseline using the same model family (\Cref{T:full-paper}). The baseline's mean reported API input
is 6,540 tokens per article, while the 40-article accounting table concerns a different sample
and an estimated rather than billed input measure. We do not combine those two measures into a
sample-specific cost ratio. The baseline also lacks a robust chunking/retry protocol; two replies
were recorded as empty or unparseable. Consequently, the comparison includes both semantic
extraction and response-format robustness and is not evidence that a well-engineered full-article
LLM baseline, or a stronger model, would have the same performance.
